\section{Methods}
\label{sec:methods}

This section describes the methodological pipeline for constructing a graph-based 
species distribution model (SOM-GNN-SDM) for Swiss flora at 30\,m resolution, 
following the framework proposed by \citet{Wu2025}. The pipeline consists of three 
stages: 
landscape patch construction via Self-Organising Map clustering and connected-
component analysis, (2)~a Random Forest baseline, and (3)~the SOM-GNN-SDM itself.

\subsection{Landscape Patch Construction}
\label{sec:patches}

The landscape patch construction transforms the continuous 30\,m raster feature 
stack into a graph representation where nodes correspond to ecologically 
homogeneous habitat patches and edges encode spatial adjacency.

\subsubsection{SOM Clustering}
\label{sec:som}

A Self-Organising Map \citep[SOM;][]{Kohonen1982} was trained to identify 
recurring landscape types from the 12 selected environmental features. The SOM 
was configured as a $64 \times 64$ hexagonal, toroidal grid and initialised with 
PCA. Training used a subsample of 500\,000 pixels (drawn from the ${\sim}160$\,M
valid pixels in the study extent), preprocessed with a \texttt{RobustScaler} and 
domain-specific post-scaling (land-cover fractions up-weighted $3\times$; 
forest-edge distance clipped to $[-5, 5]$).

The SOM was trained for 150 epochs with a neighbourhood radius decaying from 64 
to 16 and a learning rate decaying from 0.05 to 0.016. Convergence was assessed 
on a held-out 20\% validation set using quantisation error (QE) and topographic 
error (TE).

After training, the codebook vectors were hierarchically clustered using Ward's 
linkage, and the dendrogram was cut to yield 30 landscape types. Every valid 
pixel in the study extent was then assigned to its best-matching unit (BMU) via 
a FAISS exact-$L_2$ index, producing a categorical landscape-type raster at 30\,
m resolution.

\subsubsection{Connected Components and Patch Merging}
\label{sec:components}

Spatially contiguous regions of the same landscape type were identified using 
connected-component labelling (\texttt{skimage.measure.label}). To avoid an 
excessive number of tiny fragments, patches smaller than 100 pixels 
(${\sim}0.09$\,km$^2$) were merged into their most similar neighbour using a 
mapping-table approach. This procedure yielded \textbf{166,464 landscape patches} 
covering the full study extent.

For each patch, the mean value of the 12 selected environmental features was 
computed from the underlying raster stack, producing a $166{,}464 \times 12$ 
node-feature matrix.

\subsubsection{Patch Network}
\label{sec:network}

An undirected graph was constructed by connecting patches that share at least 
one boundary pixel. The resulting \texttt{NetworkX} graph contains 
\textbf{166,464 nodes} and \textbf{465,833 edges}. Node attributes are the mean 
environmental features; no edge weights were used.

\subsection{Baseline: Random Forest}
\label{sec:rf}

A per-species Random Forest \citep{Breiman2001} classifier served
as the non-spatial baseline. For each of the 2,696 qualifying
species ($\geq 20$ presence patches), the following procedure was
applied:

\begin{enumerate}
    \item \textbf{Presence labels.} GBIF occurrence coordinates were mapped to 
    the patch-label raster, yielding a set of presence patches per species.
    \item \textbf{Pseudo-absence generation.} Background patches were sampled at 
    a 3:1 ratio (absence\,:\,presence) from patches without recorded occurrences, 
    using a fixed random seed for reproducibility.
    \item \textbf{Model training.} A \texttt{scikit-learn} Random Forest 
    was trained on the patch-level mean features using tuned hyperparameters 
    (see below).
    \item \textbf{Evaluation.} Performance was assessed via 5-fold cross-validated 
    metrics: AUC, accuracy, precision, recall, F1, MCC, Cohen's $\kappa$, and TSS.
\end{enumerate}

Hyperparameter tuning was performed via \texttt{GridSearchCV} on the same 
six representative species used for the GNN architecture search, covering 
diverse habitat types across Switzerland. The best hyperparameters from 
this limited tuning were then applied uniformly to all
2,696 species.

\subsection{GNN-SDM}
\label{sec:gnn}

\subsubsection{Model Architecture}
\label{sec:architecture}

The SOM-GNN-SDM is based on GraphSAGE \citep{Hamilton2017}, implemented in 
PyTorch Geometric \citep{Fey2019}. At each layer $l$, a node $v$ aggregates 
information from its neighbourhood $\mathcal{N}(v)$ using the mean aggregator 
and combines it with its own representation:
\begin{align}
    \mathbf{h}^{(l)}_{\mathcal{N}(v)} &= \frac{1}{|\mathcal{N}(v)|} 
        \sum_{u \in \mathcal{N}(v)} \mathbf{h}^{(l-1)}_u 
        \label{eq:sage-agg} \\
    \mathbf{h}^{(l)}_v &= \sigma\!\left( 
        \mathbf{W}^{(l)} \cdot 
        \mathrm{CONCAT}\!\left( 
            \mathbf{h}^{(l-1)}_v,\; 
            \mathbf{h}^{(l)}_{\mathcal{N}(v)} 
        \right) 
    \right) \label{eq:sage-update}
\end{align}
where $\mathbf{h}^{(0)}_v = \mathbf{x}_v$ is the input feature vector of 
node $v$, $\mathbf{W}^{(l)}$ is a learnable weight matrix, and $\sigma$ is 
the LeakyReLU activation function. Dropout ($p = 0.2$) is applied after each 
layer during training. The final node embeddings are passed through a linear 
layer with sigmoid activation to produce a habitat-suitability score 
$\hat{y}_v \in [0, 1]$ for every patch:
\begin{equation}
    \hat{y}_v = \mathrm{sigmoid}\!\left( 
        \mathbf{w}^\top \mathbf{h}^{(L)}_v + b 
    \right) \label{eq:sage-output}
\end{equation}
where $L$ is the number of GraphSAGE layers.

An architecture search was conducted on six representative species spanning 
diverse habitats (\textit{Picea abies}, \textit{Fagus sylvatica}, 
\textit{Trifolium pratense}, \textit{Potentilla erecta}, 
\textit{Phragmites australis}, \textit{Senecio inaequidens}). 
Six architectures were compared:

\begin{itemize}
    \item 1-layer: $[16]$ (417 parameters)
    \item 2-layer: $[24, 12]$ (1,201 parameters)
    \item Paper default: $[24, 18, 8]$ (1,787 parameters)
    \item Wider: $[48, 32, 16]$ (5,361 parameters)
    \item Extra-wide: $[64, 48, 32]$ (10,929 parameters)
    \item 4-layer: $[32, 24, 16, 8]$ (3,417 parameters)
\end{itemize}

The extra-wide 3-layer architecture $[64, 48, 32]$ consistently achieved the 
highest mean AUC across all test species and was selected for production training. 
This is notably wider than the $[24, 18, 8]$ architecture proposed by \citet{Wu2025} 
for virtual fauna species, suggesting that flora---with their sessile nature and 
finer-grained habitat associations---benefit from greater representational 
capacity per layer rather than additional depth.

\subsubsection{Training Procedure}
\label{sec:training}

For each species, the GNN-SDM was trained independently using the following pipeline:

\begin{enumerate}
    \item \textbf{Background selection via random pseudo-absence.}
    Non-presence patches were sampled at a 3:1 ratio (background\,:\,presence) 
    using a fixed random seed. An initial design based on One-Class SVM background 
    selection \citep[as in][]{Wu2025} was evaluated during development but 
    abandoned in favour of random pseudo-absence for the production run. The 
    random strategy ensures a fair comparison with the RF baseline (which uses 
    same pseudo-absence protocol) and proved more robust across the full species 
    set, avoiding the computational overhead and occasional instability of 
    fitting a kernel SVM on high-dimensional patch features for
    each of 2,696 species.

    \item \textbf{PageRank weighting.}
    Presence patches were weighted by their PageRank centrality in the landscape 
    graph, following \citet{Wu2025}. This assigns higher importance to 
    structurally well-connected patches, reflecting their greater ecological 
    significance as habitat corridors or core areas. Background patches received 
    uniform weight of 1.0.

    \item \textbf{Train/validation split.}
    The labelled patches (presence $+$ background) were shuffled and split 80/20 
    into training and validation sets.

    \item \textbf{Optimisation.}
    The model was trained with the Adam optimiser (learning rate $= 0.001$) 
    using a weighted mean-squared-error loss:
    \[
        \mathcal{L} = \frac{1}{|\mathcal{T}|} \sum_{i \in \mathcal{T}} w_i \left( \hat{y}_i - y_i \right)^2
    \]
    where $\mathcal{T}$ is the training set, $w_i$ the PageRank-derived weight, 
    $\hat{y}_i$ the predicted suitability, and $y_i \in \{0, 1\}$ the label.

    \item \textbf{Early stopping.}
    Validation AUC was evaluated every 10 epochs. Training was halted if no 
    improvement occurred for 50 consecutive epochs (patience). The model state 
    with the best validation AUC was restored for prediction. Maximum training 
    duration was 500 epochs.

    \item \textbf{Prediction.}
    The best model was used to predict habitat suitability for all 166,464 
    patches, yielding a continuous suitability surface per species.
\end{enumerate}

\subsubsection{Evaluation}
\label{sec:evaluation}

To evaluate the GNN-SDM on the same footing as the RF baseline, a post-hoc 
evaluation was performed for each species: presence patches and a 3:1 random 
sample of non-presence patches were combined, and the predicted suitability 
scores were evaluated using the same eight metrics (AUC, accuracy, precision, 
recall, F1, MCC, Cohen's $\kappa$, TSS) with a classification threshold of 0.5.

\subsubsection{Computational Setup}
\label{sec:compute}

Production training of all 2,696 species was executed as an AWS
SageMaker Training Job on a GPU instance (NVIDIA GPU,
CUDA-accelerated). The job supported checkpointing for
spot-instance interruptions and completed in approximately 5.7
hours of GPU time (mean 7.6\,s per species, 231 epochs on
average). Node features were standardised (zero mean, unit variance) 
before training, and the full graph (166,464 nodes, 931,666 directed edges) was 
loaded onto the GPU for efficient message passing.

\subsection{Conservation Prioritisation}
\label{sec:conservation}

The final stage of the pipeline combines the GNN-SDM habitat suitability 
predictions with circuit-theory connectivity analysis to produce a multi-species 
conservation priority surface, which is then overlaid with existing protected 
areas to identify protection gaps.

\subsubsection{Complementarity of Suitability and Connectivity}
\label{sec:complementarity}

One might ask whether the circuit-theory step is redundant 
given that GraphSAGE already incorporates neighbourhood information during 
message passing. The two components are, however, fundamentally complementary:

\begin{itemize}
    \item \textbf{GNN suitability} is a \emph{local} property. GraphSAGE 
    aggregates features from the 1--2 hop neighbourhood to predict whether a 
    given patch is suitable habitat for a species. The output answers: 
    ``Is this place good habitat?''

    \item \textbf{Circuit-theory current flow} is a \emph{global network} 
    property. It treats the entire landscape graph as an electrical circuit 
    and computes which patches act as irreplaceable bridges between isolated 
    population clusters. The output answers: ``Is this place critical for 
    maintaining connectivity between populations?''
\end{itemize}

These two perspectives are orthogonal. A patch may have moderate suitability 
(GNN score ${\sim}0.3$) yet carry very high current because it lies in the only 
traversable corridor between two populations. Conversely, a patch in the core 
of a large population may have high suitability but low current because 
alternative paths exist. The highest conservation value lies a patches that score
highly on both dimensions.

\subsubsection{Circuit-Theory Connectivity}
\label{sec:circuit}

For each species, the landscape graph is treated as a resistor network where 
the conductance of each edge $(u, v)$ is defined as the mean suitability of 
the two adjacent patches:
\begin{equation}
    g_{uv} = \max\!\left( 10^{-4},\; 
        \frac{\hat{y}_u + \hat{y}_v}{2} 
    \right)
    \label{eq:conductance}
\end{equation}
This ensures that current preferentially flows through high-suitability regions 
(low resistance), reflecting the ecological principle that organisms move more 
easily through suitable habitat.

Source and ground nodes are defined by the species' actual observation patches. 
These patches are grouped into connected components in the landscape graph; 
components with $\geq 3$ patches are ranked by size, and current flow is solved 
between the top $k = 6$ component pairs using Kirchhoff's circuit laws. 
Specifically, for each source--ground pair, one unit of current is injected 
uniformly across source nodes and extracted uniformly at ground nodes. The 
resulting system of linear equations (conductance Laplacian) is solved for node 
voltages using a sparse direct solver (\texttt{scipy.sparse.linalg.spsolve}), 
and the absolute current through each node is accumulated across all pairs.

Species with fewer than 2 qualifying components are excluded, as connectivity 
cannot be meaningfully assessed.

\subsubsection{Species Selection and Filtering}
\label{sec:species-selection}

The gap analysis targets nationally threatened species according
to the Swiss Red Lists of vascular plants
\citep{InfoFlora2016} and bryophytes \citep{InfoFlora2022},
published by Info Flora on behalf of BAFU (Federal Office for
the Environment). These lists apply IUCN criteria at the
\emph{national} level, as a species can be globally Least Concern
yet nationally Critically Endangered in Switzerland. 

Of the 2,696 GNN-trained species, 489 carry a nationally
threatened status (CR, EN, VU, or NT). Two additional filters
are applied:

\begin{enumerate}
    \item \textbf{Minimum spatial structure.} Only species with
    $\geq 50$ presence patches are retained, ensuring sufficient
    spatial extent for meaningful circuit-theory analysis. This
    yields \textbf{285 species} (CR\,=\,2, EN\,=\,20, VU\,=\,93,
    NT\,=\,170).

    \item \textbf{Minimum connectivity.} During computation,
    species whose observation patches form fewer than 2 connected
    components of size $\geq 3$ are excluded, as the
    circuit-theory framework requires at least one
    source--ground pair. This removes a further 44 species,
    leaving 241 species with non-zero current flow.
\end{enumerate}

\subsubsection{Normalisation and Aggregation}
\label{sec:aggregation}

Per-species current flow values span vastly different scales
depending on graph structure and population fragmentation. To
ensure each species contributes on a comparable scale to the
multi-species aggregation, each species' current vector is
mean-normalised:
\begin{equation}
    \tilde{c}_{v} = \frac{c_v}{\bar{c}}
    \label{eq:normalisation}
\end{equation}
where $\bar{c} = \frac{1}{N}\sum_{v} c_v$ is the species-level
mean current. Values are thus expressed as multiples of the
species mean: $\tilde{c}_v = 1$ indicates average flow, while
$\tilde{c}_v \gg 1$ identifies critical pinch-points. This
allows comparison across species whose absolute current
magnitudes differ: a species with many observed
patches generates higher raw current than one with few
observations, yet both may contain equally important
bottlenecks. Mean normalisation preserves the relative
importance of pinch-points within each species while placing
all species on a common scale. Species whose mean current is
effectively zero ($\bar{c} < 10^{-12}$) are excluded, leaving
241 of 285 species with non-zero flow.

Species are weighted by national Red List category
(CR\,=\,4, EN\,=\,3, VU\,=\,2, NT\,=\,1) to reflect
conservation urgency. The multi-species priority surface
is then computed as:
\begin{equation}
    P_v = \frac{1}{2}\left( 
        \frac{\sum_s w_s \hat{y}_{v,s}}{\sum_s w_s} + 
        \frac{\sum_s w_s \tilde{c}_{v,s}}{\sum_s w_s} 
    \right)
    \label{eq:priority}
\end{equation}
where $w_s$ is the threat-category weight of species $s$,
$\hat{y}_{v,s}$ the GNN suitability, and $\tilde{c}_{v,s}$
the normalised current flow at patch $v$.

\subsubsection{Protected Area Overlay}
\label{sec:overlay}

The Swiss protected area network (WDPA; 10,522 polygons) is rasterised onto 
the 30\,m patch grid. Each patch is assigned the fraction of its constituent 
pixels that fall within any protected area polygon. The analysis is restricted 
to patches within Swiss territory (majority of pixels inside the national 
boundary) to avoid false negatives from neighbouring countries.

Conservation gaps are quantified by examining what fraction of high-priority 
patches (top 5\% and top 10\% of the priority surface) currently lack any 
protected-area coverage. These unprotected high-priority patches represent 
the most urgent candidates for protection expansion.
